
En
appliquant 3.7, on trouve :
@@M0@@
La fonctorialité de @@M1@@ et de @@M2@@ entraîne :
@@M3@@
D'après 2.7, on a :
@@M4@@
La multiplicativité du caractère de Todd entraîne donc :
@@M5@@
On obtient donc finalement :
@@M6@@
ce qui est le résultat annoncé.

\textbf{3.9.} Si on applique 3.8 à la factorisation de @@M7@@ donnée en 3.1, on voit qu'il suffit, pour démontrer le théorème de Riemann--Roch, de le démontrer dans le cas d'une immersion régulière de codimension constante, et dans le cas du morphisme structural : @@M8@@.

En outre, on peut supposer que l'immersion @@M9@@ satisfait la condition de VII 4.1. Soit en effet @@M10@@ la section de @@M11@@ introduite en VII 4.2. Alors @@M12@@ est une immersion régulière vérifiant la condition de VII 4.1 (VII 4.2) et d'après 3.8, il suffit de montrer 3.6 pour @@M13@@, @@M14@@ et @@M15@@.

On va donc prouver 3.6 en prouvant les deux cas particuliers suivants :
\begin{enumerate}[label=\roman*)]
\item @@M16@@ est une immersion régulière de codimension constante, satisfaisant la condition de VII 4.1;
